\documentclass{standalone}
\usepackage{circuitikz}

\begin{document}

\begin{circuitikz}

% Línea principal con resistencias
\draw
(0,0) node[left]{$T_j$}
to[short, o-] (1,0)
to[R] (3,0)
to[R] (5,0)
to[R] (7,0)
to[R] (9,0)
to[short] (9,0) node[right]{$T_{case}$};

% Capacitores en paralelo (ramas independientes)
\draw (1,-2) to[C] (3,-2);
\draw (3,-2) to[C] (5,-2);
\draw (5,-2) to[C] (7,-2);
\draw (7,-2) to[C] (9,-2);

% Conexiones entre componentes
\draw (1,0) to[short] (1,-2);
\draw (3,0) to[short] (3,-2);
\draw (5,0) to[short] (5,-2);
\draw (7,0) to[short] (7,-2);
\draw (9,0) to[short] (9,-2);

% Conexión a T_case
\draw (9,-2) to[short] (9,-2.5)
node[ground]{}
node[below]{};

\end{circuitikz}

\end{document}